% Section 2: Related Work (~1 page, ~750 words)
% Outline only — prose pass deferred. Three subsections, ~250 words each.

\section{Related Work}
\label{sec:related_work}

\subsection{Ensemble methods in tabular reinforcement learning}
\label{sec:rw_tabular}

\guideblock{One paragraph. Walk through the tabular ensemble line:
\citet{wiering2008} (first systematic study of voting rules in
$Q$-learning), \citet{fausser2011} (ensemble methods with function
approximation), \citet{lan2020maxmin} (maxmin $Q$-learning as
ensemble-style estimation-bias control), \citet{peer2021ebql} (ensemble
bootstrapping for $Q$-learning). The distinction this paper draws: prior
tabular ensemble work either lives in toy gridworld / small MDP regimes
\emph{or} uses algorithmic diversity (different update rules per member);
we use seed-only diversity in a non-toy stochastic two-player game.}

\subsection{Ensemble methods in deep reinforcement learning}
\label{sec:rw_deep}

\guideblock{One paragraph. Cover the deep-RL ensemble line:
\citet{osband2016bootstrapped} (bootstrapped DQN, posterior-sampling
exploration), \citet{anschel2017avgdqn} (averaged DQN for variance
reduction), \citet{chen2021redq} (randomized ensembled double
$Q$-learning), \citet{lee2021sunrise} (unified ensemble framework).
Close with the \citet{lin2024curse} ``Curse of Diversity'' failure mode:
deep ensembles sharing a replay buffer can underperform their own
single-agent baseline. Set up the contrast: our independent-table setup
has no shared buffer and is shown to be immune by construction (full
discussion in Section~\ref{sec:d_curse}).}

\subsection{Reinforcement learning for Pokemon}
\label{sec:rw_pokemon}

\guideblock{One paragraph in roughly chronological order:
\citet{lee2017showdown} (early Showdown AI competition),
\citet{kalose2018} (Stanford project, early Pokemon RL writeup),
\citet{sahovic2020pokeenv} (\texttt{poke-env}, the library used here),
\citet{simoes2020} (deep RL on the Showdown simulator),
\citet{wang2024mit} (MIT M.Eng.\ thesis on random-battle RL),
\citet{hu2024pokellmon} (LLM-as-policy approach),
\citet{karten2025pokechamp} (expert-level minimax LM agent),
\citet{grigsby2025metamon} (transformer offline RL),
\citet{vgcbench2025} (VGC team-strategy benchmark). Close with the gap
this paper addresses: to our knowledge, no prior paper trains $K$
independent tabular $Q$-function agents under identical hyperparameters
and combines them by inference-time voting on Pokemon battles.}
